\documentclass[10pt]{article}
\providecommand{\FIGDIR}{fig}
\newcommand{\DIGESTDATE}{15 August 2026}
\newcommand{\PERIODSTART}{9 August 2026}
\newcommand{\PERIODEND}{15 August 2026}
% ===================== PM Research Watch - shared preamble =====================
\usepackage[T1]{fontenc}
\usepackage[utf8]{inputenc}
\usepackage{amsmath,amssymb}
\usepackage{tgpagella}
\usepackage{tgheros}
\renewcommand{\sfdefault}{qhv}
\renewcommand{\ttdefault}{lmtt}
\usepackage{microtype}
\usepackage[english]{babel}

\usepackage[a4paper,top=17mm,bottom=17mm,left=17mm,right=17mm,
            headsep=6mm,footskip=10mm]{geometry}
\usepackage{graphicx}
\usepackage{xcolor}
\usepackage{colortbl}
\usepackage{tikz}
\usetikzlibrary{calc}
\usepackage[most]{tcolorbox}
\usepackage{titlesec}
\usepackage{fancyhdr}
\usepackage{booktabs}
\usepackage{array}
\usepackage{longtable}
\usepackage{enumitem}
\usepackage{multicol}
\usepackage{ragged2e}
\usepackage[hidelinks]{hyperref}

% ------------------------------------------------------------------ palette
\definecolor{Ink}    {HTML}{14262E}
\definecolor{Deep}   {HTML}{0F4C5C}
\definecolor{Teal}   {HTML}{1F7A8C}
\definecolor{Sky}    {HTML}{4FA3C4}
\definecolor{Sage}   {HTML}{6FA287}
\definecolor{Amber}  {HTML}{E3A008}
\definecolor{Clay}   {HTML}{B5643C}
\definecolor{Coral}  {HTML}{D1495B}
\definecolor{Violet} {HTML}{6A5B9E}
\definecolor{Slate}  {HTML}{4A6572}
\definecolor{Paper}  {HTML}{FBF9F5}
\definecolor{Mist}   {HTML}{EEEAE2}
\definecolor{Rule}   {HTML}{D9D3C9}

\pagecolor{Paper}
\color{Ink}

\hypersetup{colorlinks=true, linkcolor=Deep, urlcolor=Teal, citecolor=Teal,
            pdftitle={PM Research Watch}, pdfauthor={Automated daily digest}}

% ------------------------------------------------------------------ headings
\titleformat{\section}
  {\sffamily\bfseries\large\color{Deep}}{}{0pt}
  {}[{\vspace{1.2mm}\color{Amber}\titlerule[1.1pt]}]
\titlespacing*{\section}{0pt}{5.5mm}{2.8mm}

\titleformat{\subsection}
  {\sffamily\bfseries\normalsize\color{Teal}}{}{0pt}{}
\titlespacing*{\subsection}{0pt}{3.6mm}{1.4mm}

\setlist[itemize]{leftmargin=4.2mm, itemsep=1.1mm, topsep=1.3mm, parsep=0pt,
                  label=\textcolor{Teal}{\footnotesize$\blacktriangleright$}}

% ------------------------------------------------------------------ header/footer
\pagestyle{fancy}
\fancyhf{}
\renewcommand{\headrulewidth}{0pt}
\renewcommand{\footrulewidth}{0pt}
\fancyhead[L]{\sffamily\scriptsize\color{Slate}PM RESEARCH WATCH\;\textcolor{Rule}{\textbar}\;\DIGESTDATE}
\fancyhead[R]{\sffamily\scriptsize\color{Slate}Particulate matter monitoring \& health effects}
\fancyfoot[C]{\sffamily\scriptsize\color{Slate}\thepage}

% ------------------------------------------------------------------ components
% metric tile
\newcommand{\metric}[3]{%
  \begin{tcolorbox}[enhanced, colback=white, colframe=Rule, boxrule=0.5pt,
    arc=1.6mm, left=2.2mm, right=2.2mm, top=1.8mm, bottom=1.8mm,
    borderline west={1.7mm}{0pt}{#3}, nobeforeafter, valign=center,
    height=17mm, width=\linewidth]
    \RaggedRight
    {\sffamily\bfseries\fontsize{17}{17}\selectfont\color{#3}#1}\\[0.5mm]
    {\sffamily\fontsize{6.5}{7.4}\selectfont\color{Slate}#2}
  \end{tcolorbox}}

% pull-quote / key insight box
\newtcolorbox{keybox}[1][]{enhanced, colback=Mist, colframe=Mist,
  arc=1.6mm, left=3.2mm, right=3.2mm, top=2.4mm, bottom=2.4mm,
  borderline west={1.7mm}{0pt}{Amber}, #1}

% subtopic band
\newcommand{\band}[2]{%
  \vspace{3.2mm}\noindent
  \begin{tcolorbox}[enhanced, colback=#2, colframe=#2, arc=1.2mm,
    left=2.6mm, right=2.6mm, top=1.1mm, bottom=1.1mm, boxrule=0pt,
    width=\linewidth, nobeforeafter]
    {\sffamily\bfseries\small\color{white}#1}
  \end{tcolorbox}\vspace{1.6mm}}

% one paper entry: {short cite}{title}{journal}{doi}{body}
\newcommand{\paperentry}[5]{%
  \par\vspace{1.9mm}\noindent
  {\sffamily\bfseries\footnotesize\color{Ink}#1}\;%
  {\sffamily\fontsize{7}{8}\selectfont\color{Slate}#3}\\[0.4mm]
  {\footnotesize\itshape\color{Deep}#2}\\[0.6mm]
  {\footnotesize #4}\\[0.3mm]
  {\fontsize{6.6}{7.6}\selectfont\color{Slate}\texttt{doi:}\,\href{https://doi.org/#5}{\texttt{\detokenize{#5}}}}
  \par}

% caption style
\newcommand{\figcap}[1]{\par\vspace{1.0mm}
  {\fontsize{7.2}{8.6}\selectfont\color{Slate}\textsf{\textbf{Fig.}}\; #1\par}\vspace{1.5mm}}

\providecommand{\permil}{\textperthousand}

\setlength{\parindent}{0pt}
\setlength{\parskip}{0pt}
\renewcommand{\arraystretch}{1.18}


\begin{document}

% ==================================================================== MASTHEAD
\thispagestyle{empty}
\begin{tikzpicture}[remember picture, overlay]
  \fill[Deep] (current page.north west) rectangle
      ($(current page.north east)+(0,-67mm)$);
  \fill[Teal] ($(current page.north west)+(0,-67mm)$) rectangle
      ($(current page.north east)+(0,-68.4mm)$);
  \foreach \x/\y/\r in {22/12/0.9, 41/26/1.5, 63/9/0.7, 88/31/1.1, 112/17/1.9,
                        138/29/0.8, 159/13/1.3, 178/54/1.0, 196/21/1.6,
                        30/44/0.6, 74/57/1.2, 125/48/0.9, 168/6/0.8, 191/59/0.7,
                        52/18/1.0, 100/60/0.6, 148/20/1.1, 66/49/1.4, 205/45/0.8}
    \fill[Sky, opacity=0.30] ($(current page.north west)+(\x mm,-\y mm)$) circle (\r mm);
\end{tikzpicture}

\vspace*{4mm}
{\sffamily\fontsize{7.6}{9}\selectfont\color{Sky}%
 WEEKLY ROLLUP\;\textbullet\;PhD RESEARCH WATCH\par}
\vspace{2.2mm}
{\sffamily\bfseries\fontsize{27}{29}\selectfont\color{white}%
 Particulate Matter\par}
{\sffamily\fontsize{27}{29}\selectfont\color{Sky}%
 Week in Review\par}
\vspace{5.0mm}
{\sffamily\fontsize{8.4}{10}\selectfont\color{white}%
 Week ending \DIGESTDATE\;\;\textcolor{Teal}{\textbar}\;\;\PERIODSTART{} to \PERIODEND\;\;%
 \textcolor{Teal}{\textbar}\;\;134 records pooled from 7 daily issues\par}
\vspace{18mm}

% -------------------------------------------------------------- metric strip
\noindent
\begin{minipage}[t]{0.190\linewidth}\metric{134}{RECORDS POOLED\\(134 UNIQUE DOIs)}{Deep}\end{minipage}\hfill
\begin{minipage}[t]{0.190\linewidth}\metric{45\%}{ON THE MEASUREMENT\\SIDE (60 RECORDS)}{Teal}\end{minipage}\hfill
\begin{minipage}[t]{0.190\linewidth}\metric{10}{SUBTOPICS\\ACTIVE}{Sky}\end{minipage}\hfill
\begin{minipage}[t]{0.190\linewidth}\metric{32}{POOLED EFFECT\\ESTIMATES}{Amber}\end{minipage}\hfill
\begin{minipage}[t]{0.190\linewidth}\metric{0}{NEW TRIAL\\REGISTRATIONS}{Clay}\end{minipage}

\vspace{6mm}

% -------------------------------------------------------------- week in one box
\section*{\sffamily\bfseries\color{Deep}The week in one paragraph}
\vspace{-1.5mm}{\color{Teal}\rule{\linewidth}{1.1pt}}\vspace{2.5mm}

\begin{keybox}
{\sffamily\bfseries\small\color{Deep}The measurement side gave ground this week --- 45\%
against 54\% --- and what replaced it was clinical-endpoint work in organ systems the
watch had not seen: kidney, thyroid, brain white matter, Parkinson's. The instrument
literature meanwhile turned from where a calibration fails to \emph{how long} a sensor
survives, which is the same question one level down.}\\[1.4mm]
\footnotesize
Last week's finding was that four groups independently stopped publishing new low-cost
sensor correction equations and started publishing the failure boundaries of existing
ones. This week that inquiry got more specific and more uncomfortable.
\textbf{Benjamin et al.} (\textit{Sensors}) show that a low-cost CO and PM$_{2.5}$ unit
calibrated in one environmental regime does not extrapolate to another --- the failure is
in the environment, not the reference. \textbf{Balasubramanian et al.} fit a
PurpleAir correction across the whole United States, which is the opposite move and is
useful precisely as a contrast: a national model is a statement that regional structure
can be absorbed, and Benjamin's result is a statement that it cannot.
\textbf{Fawaz et al.} address a failure mode the correction literature does not touch at
all --- surface fouling of a SAW microparticle sensor in a cascade impactor, cleaned in
situ by SAW-induced droplet displacement --- and \textbf{Kirchhoff et al.} re-examine the
working fluid in condensation particle counters. Two of those four are about the sensor
surviving contact with real air, not about fitting its output.

On the health side the week's five tier-A clinical records land in five different organ
systems: \textbf{Mandal} (eGFR decline, 12{,}271 Indian adults), \textbf{Wang et al.}
(thyroid, $n=24{,}333$), \textbf{Sullivan} (childhood white-matter microstructure),
\textbf{Delaney} (first Parkinson's hospitalisation) and \textbf{Chow} (blood metabolomic
signatures linking air pollution to lung cancer). \textbf{Not one of the five measured the
exposure it analysed.} All five inherit a modelled surface --- satellite, land-use
regression or a spatiotemporal product --- and none reports the model's error at the
concentrations where the effect is claimed. \textbf{Across 134 records and 34 sensing
papers, the count of records that validated an instrument against a health endpoint is
again zero}, for the third consecutive weekly rollup.
\end{keybox}

\vspace{3mm}

% -------------------------------------------------------------- trend & inflection
\section*{\sffamily\bfseries\color{Deep}Subtopic trend and inflection}
\vspace{-1.5mm}{\color{Teal}\rule{\linewidth}{1.1pt}}\vspace{2.5mm}

\noindent\includegraphics[width=\linewidth]{\FIGDIR/w1_subtopic_trend.png}
\figcap{Per-issue subtopic counts across the week from \texttt{state/metrics.csv}.
\textbf{Read the denominators first}: daily totals were 12, 16, 26, 28, 17, 14 and 21 ---
a factor of 2.3 between the largest and smallest day, and a much flatter week than
2--8 August, whose range was 10 to 31. The 11--12 August peak (26 and 28) is an Elsevier
\textit{Atmospheric Environment} deposit batch, visible in the register as nine
\texttt{atmosenv.2026.1222xx/1223xx} DOIs arriving across two days, most without
abstracts. \textbf{That is a publisher's release schedule, not the field.} The 9 August
floor of 12 is a Sunday. Once those two artefacts are removed the week has no inflection
at all, and the honest summary is that subtopic composition was stable and the total
declined slightly.}

\footnotesize
\textbf{Fading, and the one number worth carrying forward.} Sensing dropped
\textbf{47 $\rightarrow$ 34} and exposure modelling \textbf{32 $\rightarrow$ 26}; the
measurement-side share fell from \textbf{54\% to 45\%}. Both of last week's peaks were
partly this pipeline's own doing --- the Crossref-by-ISSN leg was widened on 2 August and
immediately surfaced ten calibration papers --- so the fall is at least partly a return to
baseline rather than a decline in the field. \textbf{Do not read $-$13 as a signal.}

\vspace{1.2mm}
\textbf{Emerging.} \emph{Clinical endpoints outside the cardiorespiratory core.}
\textit{Other clinical endpoints} rose \textbf{6 $\rightarrow$ 11} and \textit{Neuro /
mental health} \textbf{2 $\rightarrow$ 5}, and the content is genuinely new for this
corpus: renal function (Mandal; Tan et al.\ on eGFR with nutritional and activity
indices), thyroid (Wang et al., $n=24{,}333$; a nodular-thyroid record), haematological
malignancy and CHIP, and Alzheimer's neuropathology. \textbf{This is the first week in
which the watch's clinical tail was longer than its cardiorespiratory trunk}, and it is
the shift most likely to persist, because it tracks the arrival of large linked-registry
cohorts rather than a publication batch.

\vspace{1.2mm}
\textbf{Flat but load-bearing.} Mechanistic toxicology held at exactly \textbf{11} for the
second week, occupational and indoor at \textbf{9}, cardiovascular at \textbf{7}. Three
clusters unchanged to the record across a fourteen-day window is worth stating: this
corpus has a stable substrate and a volatile surface, and almost all week-to-week movement
is in sensing, exposure modelling and burden --- the three clusters most sensitive to which
harvesting leg happened to fire.

%% LAYOUT RULE (2026-08-02): architecture donut in the whitespace here, not on its own page.
\vspace{3mm}
\noindent\includegraphics[width=\linewidth]{\FIGDIR/f2a_architecture.png}
\figcap{Pooled study architecture over 134 records. Measurement campaigns (26) and
modelling/inventory (21) lead, as in every week of this watch. \textbf{Two slices should
be read as defects rather than designs.} \textit{Metadata only} is \textbf{15 records}
--- 11\% of the week --- surfaced with no abstract in Crossref, Europe PMC or Semantic
Scholar and therefore classified but unread; nine of the fifteen are the
\textit{Atmospheric Environment} batch above. \textit{Review / synthesis} is 16 and is a
publication-cycle artefact, not new evidence. \textbf{The \textit{Trial / intervention}
slice is 2 records in 134} --- a protocol and a feasibility pilot, neither reporting an
outcome. \textbf{For the third week running, not one completed randomised intervention on
a PM exposure entered the corpus}, while 21 such trials sit in the registry watch below.}

\clearpage

% -------------------------------------------------------------- pooled analytics
\section*{\sffamily\bfseries\color{Deep}Pooled analytics}
\vspace{-1.5mm}{\color{Teal}\rule{\linewidth}{1.1pt}}\vspace{3mm}

\noindent\includegraphics[width=\linewidth]{\FIGDIR/f1_subtopics.png}
\figcap{Pooled subtopic distribution, \PERIODSTART{} to \PERIODEND. \textbf{Sensing (34)
and Exposure assessment \& modelling (26) account for 60 of 134 records, 45\% of the
week}, down nine points on the preceding week. All ten canonical clusters were active,
which is a property of a seven-day window rather than of any day --- \textit{Neuro / mental
health} was empty on five of the seven issues and still finished at 5. The clinical
clusters now form a genuine block rather than a tail: \textit{Other clinical endpoints}
(11), \textit{Respiratory \& allergic} (11), \textit{Occupational \& indoor} (9),
\textit{Cardiovascular \& metabolic} (7), \textit{Reproductive \& developmental} (5) and
\textit{Neuro} (5) total 48 records against 60 on the measurement side.}

\vspace{2mm}
\noindent\includegraphics[width=\linewidth]{\FIGDIR/f2b_endpoint.png}
\figcap{Pooled health endpoint by organ system. \textbf{The chart is misleading as drawn
and the reason is a label defect, not a finding.} ``No health endpoint'' shows 35, but
seven further free-text labels --- \texttt{None (exposure)} 12, \texttt{None (instrument)}
11, \texttt{None (exposure only)} 4, \texttt{None (emission)} 2, and one each of
\texttt{None (emissions)}, \texttt{None (algorithm)} and \texttt{None (in vitro OP assay)}
--- mean the same thing and are not in \texttt{ENDPOINT\_CANON}. \textbf{The true
no-endpoint count is 67 of 134, or 50\%}, which is the direct consequence of the 45\%
measurement weighting plus the 15 unread metadata-only records. Logged for repair in
\texttt{plots.py}. Of the records that do carry a human endpoint, respiratory (10) and
cardiovascular (10) lead, with oncologic at 7 --- the highest oncology count this watch has
recorded. \textbf{Every bar below 3 records is one or two studies}; do not read the tail.}

\clearpage

\noindent\includegraphics[width=\linewidth]{\FIGDIR/f5_forest.png}
\figcap{\textbf{Pooled effect estimates for the week --- a caterpillar plot, not a
meta-analysis.} No summary diamond is drawn and none should be. Pooling is not defensible
and no $I^2$ is reported, for the same four independent reasons as last week.
(i) \emph{Exposure contrasts differ}: per 1\,$\mu$g/m$^3$ (Moro, Atalell), per IQR, per
10\,$\mu$g/m$^3$, 99th-percentile-versus-median at cumulative lag (Shi), and whole-event
contrasts (Zhang's cold-spell/pollution scenarios). (ii) \emph{Estimands differ}: 14 RR,
10 OR, 4 RR-equivalent, 2 GMR and 2 adjusted OR share an axis but not an interpretation.
(iii) \emph{Records are not independent} --- Zhang contributes three nested contrasts from
one meta-analysis, Atalell two nested domain thresholds from one cohort, Negahi three
estimates two of which are the same data under different dependence assumptions.
(iv) \emph{Not all exposures are PM}. \textbf{The most informative single fact on this
plot is the width distribution, not the location}: the tightest intervals belong to
whole-of-population registry cohorts with per-unit contrasts near the null, and the widest
to small cohorts with large contrasts. That is a statement about study size, and reading
the ordering as an effect-size ranking would be an error.}

\vspace{2mm}
\noindent\includegraphics[width=\linewidth]{\FIGDIR/f4_heatmap.png}
\figcap{Pooled subtopic $\times$ study-architecture cross-tabulation over the full week.
134 records occupy \textbf{47 of the 130 available cells}. \textbf{The cells that stay
empty are the point.} At $n=134$ the \textit{Sensing} and \textit{Exposure assessment}
rows still contain \textbf{no \textit{Trial / intervention} cell and no cell coupling an
instrument to a health outcome} --- the third consecutive week this has held.
\textit{Mechanistic toxicology} remains almost entirely experimental with no bridge to
human biomarker work, unchanged since 27 July. Three weeks of accumulation and \textbf{575 records} since
26 July, and the coupling cell has never been occupied: at this point it is a property of
how the field is organised, not of sampling.}

\clearpage

% -------------------------------------------------------------- groups & citations
\section*{\sffamily\bfseries\color{Deep}Groups, venues and most-cited work}
\vspace{-1.5mm}{\color{Teal}\rule{\linewidth}{1.1pt}}\vspace{2.5mm}

\footnotesize
\textbf{Recurring author and group clusters.} The naive first-author count is
\textbf{misleading and must not be used}: the top entries are Liu (6), Zhang (6), Wang (4),
Chen (4), Li (3) and Wu (3), which are among the most common surnames in the world and
almost certainly denote distinct groups --- inspection of affiliations finds no two of the
six Zhang records sharing an institution. \textbf{One genuine repeat survives}: Cho
appears twice, on the Norwood community-science mobile campaign (15 August) and earlier in
the week, and those are the same programme. \textbf{No institution contributed more than
two records}, and there is no evidence of a single programmatic output stream inflating
apparent consistency. \textbf{Deduplication was clean this week: 134 records over 134
unique DOIs}, against 146 over 145 last week --- the cross-issue duplicate that shipped on
2 and 3 August has no counterpart in this period.

\vspace{2mm}
\textbf{A label defect repaired by this rollup.} Six records carried on 11 August were
stored with a journal name in the author field (\texttt{Atmos Environ 2026a--d},
\texttt{Aerosol Air Qual Res 2026}, \texttt{Atmos Pollut Res 2026}) because the Elsevier
and Springer deposits carried no abstract and the daily run did not resolve authorship.
A Crossref re-query during this rollup resolved all six --- Zhou, Habib, Hao, Li, Hsieh and
Yavuz --- and the corpus store has been corrected. \textbf{The 11 August PDF as shipped
still shows the journal names in its register}; the store, this rollup and every future
rollup show the authors. Recorded rather than silently fixed.

\vspace{2mm}
\textbf{Venue concentration.} \textit{Atmospheric Environment} is the largest single
contributor at \textbf{13 records}, and \textbf{that is one publisher deposit batch across
11--12 August, not a research surge} --- nine of the thirteen arrived on two consecutive
days without abstracts. After it: \textit{Environmental Monitoring and Assessment} 7,
\textit{Environmental Pollution} 6, then \textit{Journal of Hazardous Materials},
\textit{Atmospheric Chemistry and Physics}, \textit{Environmental Science \& Technology},
\textit{Journal of Thoracic Disease}, \textit{Atmospheric Measurement Techniques},
\textit{Sensors} and Research Square at 4 each. \textbf{The free-text venue defect
recorded last week persists}: \texttt{Research Square (preprint)}\,(4) and
\texttt{Res Sq (preprint)}\,(3) are the same venue under two spellings, so preprints are
actually the second-largest source of the week at 7 records.

\vspace{2mm}
\textbf{Most-cited work of the period.} \textbf{Not reportable, and ranking it would be
noise.} Every record entered a PubMed, Europe PMC or Crossref index within the last seven
days; citation counts are zero or near-zero across all 134 and carry no signal about
importance. A defensible most-cited list requires a lag of at least 18 months and will be
reported in the yearly issue, not weekly.

\vspace{4mm}

% -------------------------------------------------------------- trial watch
\section*{\sffamily\bfseries\color{Deep}Trial registry watch}
\vspace{-1.5mm}{\color{Teal}\rule{\linewidth}{1.1pt}}\vspace{2.5mm}

{\fontsize{7.4}{9}\selectfont\color{Slate}From \texttt{state/trials.json}. Daily issues
accumulate; the rollup is where they surface. \textbf{Zero registrations entered state
during 9--15 August} --- ClinicalTrials.gov was checked on each of the seven daily runs and
returned no PM or air-pollution registration created or updated inside any daily window.
The standing list of \textbf{21} is unchanged since 8 August. Note the hard limitation
recorded previously: the connector exposes \emph{no} last-update-date filter, so these are
relevance sweeps rather than true window queries, and \texttt{first\_seen} means ``first
seen by this pipeline''.\par}

\vspace{1.5mm}
\footnotesize
\textbf{A zero week is itself informative, with a caveat.} Seven consecutive days with no
new PM-exposure registration is consistent with the structural picture this watch has
built: the interventional literature in this field is small, almost entirely
filtration-based, and adds a few registrations a month rather than a few a week. But the
absence of an update-date filter means a \emph{silently updated} existing registration
would not surface, so read this as ``no new trial found'', not ``no registry activity''.

\vspace{1.5mm}
\textbf{Results now overdue --- the list has grown.} \textbf{NCT05994937} (NYU Langone,
portable air cleaner, glycaemic and cardiometabolic biomarkers, $n=150$) reached primary
completion \textbf{2026-03-16} with no results posted --- now \textbf{five months
overdue}, and note that its protocol was published this week as a daily record
(Wittkopp et al., \textit{Am Heart J}, 14 August) while the results remain unposted, which
is a reporting-order worth flagging. \textbf{NCT06534385} (Alaska Native Tribal Health
Consortium, $n=15$, 2026-04-01), \textbf{NCT07047430} (U.\ Washington, airport-adjacent
asthma, $n=60$, 2026-06-01) and \textbf{NCT07111208} (MARKOPOLO, noise and ultrafine
particles, $n=180$, 2026-07) are all past primary completion with nothing posted.
\textbf{NCT07479420} (cluster-randomised HEPA versus ionisation/plasma air cleaners,
$n=180$) is marked \texttt{COMPLETED} and is the single registration in state most likely
to produce a directly usable exposure-reduction estimate; it should be watched weekly.

\vspace{1.5mm}
\textbf{Publication-risk item.} \textbf{NCT06197477} (Hasselt, school air purifiers,
$n=900$ --- the largest trial in state) remains \texttt{UNKNOWN} past a December 2025
primary completion, with four cognitive co-primary endpoints and no stated multiplicity
control. Eight months of unknown status on the largest registered trial in the field is
the single biggest outstanding gap in the interventional evidence base.

\vspace{1.5mm}
\textbf{Adequately powered and still reading out.} \textbf{NCT06376994} (JHSPH, air
cleaner versus sham in COPD, $n=770$, 2029-08-31) and \textbf{NCT05874479} (NYU Langone,
blood pressure in NYC public housing, $n=440$, 2027-09) remain the only two registrations
that can plausibly detect a modest effect on a surrogate endpoint. Neither reads out
inside two years.

\clearpage

% -------------------------------------------------------------- gaps & direction
\section*{\sffamily\bfseries\color{Deep}Open gaps}
\vspace{-1.5mm}{\color{Teal}\rule{\linewidth}{1.1pt}}\vspace{2.5mm}

\footnotesize
\begin{itemize}
\item \textbf{Measurement--epidemiology coupling: zero for a third week.} Across 134
records and 34 sensing papers, \textbf{not one validated an instrument against a health
endpoint}. The nearest approach this week is Makacha et al.\ (14 August), which ran
minute-resolved \emph{personal} PM$_{2.5}$ monitoring on 343 women against facility
reference monitors and found ambient--personal concordance of \textbf{0.01--0.56 with
slopes below 1 in every setting} --- an exposure-misclassification measurement with no
health outcome attached. Put it beside the five tier-A clinical records above, none of
which measured its own exposure, and the shape of the gap is exact.

\item \textbf{Every large clinical result this week rests on an unvalidated exposure
surface.} Mandal (1\,km spatiotemporal model, Delhi annual means 118--130\,$\mu$g/m$^3$),
Moro (satellite model over rural Ghana), Atalell (SA2-level daily PM$_{2.5}$ at a
1\,$\mu$g/m$^3$ contrast), Sullivan and Delaney (modelled residential surfaces).
\textbf{None reports out-of-sample model error at the exposure levels where the effect is
claimed.} This is not a criticism of any one paper --- it is the field convention --- but
at per-unit odds ratios of 1.04--1.06 the convention is now the binding constraint on
inference.

\item \textbf{Ultrafine and speciated metrics: worse than last week.} Of 134 records,
\textbf{exactly one} classifies as ultrafine/nanoscale and four as size-resolved or number
distribution, against 48 PM$_{2.5}$-only and 16 joint PM$_{2.5}$/PM$_{10}$. Composition
and speciation reached 8, and optical properties 7 --- both up, and both driven by the
\textit{Atmospheric Environment} deposit batch rather than by sustained interest. The
deposition physics that determines dose remains the most under-served measurement axis in
this corpus.

\item \textbf{Geographic coverage against attributable burden: unchanged and still
inverted.} China 31, North America 20, Europe 19 against \textbf{South Asia 6 and
sub-Saharan Africa 5}. South Asia and sub-Saharan Africa together carry a far larger share
of global PM-attributable DALYs than China, North America and Europe combined, and
contributed 11 of 134 records. Two of those 11 --- Mandal and Moro --- are this week's
strongest clinical records, which makes the scarcity more striking, not less.

\item \textbf{Unread records are 11\% of the corpus.} Fifteen records carry no abstract
from any checked source and were classified but not evaluated; a Crossref re-query during
this rollup resolved \textbf{0 of 15} abstracts, though it did resolve six authorships.
The nine \textit{Atmospheric Environment} records among them are exactly the
instrumentation and source-apportionment work this watch most wants to read.
\textbf{Absence of criticism in the register does not mean a record is sound.}

\item \textbf{Negative and null results: present, which is a good sign.} Moro reports a
clean null for metabolic syndrome and all its non-hypertension components alongside a
positive hypertension estimate; Hung reports no significant mediation through any measured
inflammatory cytokine, against the default assumption in that literature; Atalell reports
null prenatal and cumulative windows alongside a positive early-childhood window. A corpus
with no nulls is a reporting-bias signal. This one has them and they are load-bearing.
\end{itemize}

\vspace{3mm}
\begin{keybox}
{\sffamily\bfseries\footnotesize\color{Deep}Where this is heading --- judgement, clearly
marked as such}\\[1.2mm]
\footnotesize
\textbf{The instrumentation literature is descending the reliability stack, and that is
the most useful thing it has done in three weeks.} Two weeks ago the question was whether
a correction fits. Last week it was where a correction stops working. This week the
questions were sensor surface fouling (Fawaz), working-fluid choice in a CPC (Kirchhoff),
environmental extrapolation limits (Benjamin) and deployment-scale correction structure
(Balasubramanian). Those are maintenance, drift and physical-integrity questions --- the
things that actually determine whether a five-year network deployment yields usable data,
and the things a held-out $R^2$ on a two-week co-location cannot address. Confidence:
moderate. Four papers over seven days is a thin base, and the same caution applies as
last week --- some of what looks like a turn is this pipeline's harvesting legs firing
differently.

\vspace{1.4mm}
\textbf{What would change the picture.} The same single paper as last week: a calibrated
low-cost network reporting an exposure--response against a clinical endpoint. Nothing in
the 21 registered trials will supply it. The likelier route is now visible in this week's
corpus --- a large linked-registry cohort (the kind producing the thyroid, renal and
Parkinson's results) retrofitting a dense sensor layer onto outcome ascertainment it
already has. That requires no new trial, only a cohort willing to instrument its
participants' addresses.

\vspace{1.4mm}
\textbf{What a distributed low-cost PM sensing programme should do differently.} Three
things this week's corpus argues for, and they are cheaper than the alternative.
\textbf{(1)} Budget for \emph{maintenance} as a measured quantity, not an operational
afterthought: Fawaz's in-situ cleaning and Benjamin's extrapolation limits both say node
performance degrades in ways a deployment-time calibration does not capture, so schedule
re-co-location and report performance as a function of time since it. \textbf{(2)} Report
exposure-surface error at the concentrations of interest, because every large clinical
result this week would have been strengthened by one number nobody supplied. \textbf{(3)}
Site the next network where the burden is, not where the reference monitors are --- South
Asia and sub-Saharan Africa produced this week's two best clinical results on the worst
exposure data in the corpus, and that is the highest-leverage measurement gap available.
\textbf{The week was quieter than the last two, and more clinically diverse than any week
so far.}
\end{keybox}

\clearpage

% ==================================================================== DIGEST
\section*{\sffamily\bfseries\color{Deep}Paper-level digest --- the week by subtopic}
\vspace{-1.5mm}{\color{Teal}\rule{\linewidth}{1.1pt}}\vspace{2mm}

\footnotesize
Every record of the period, grouped by subtopic. \textbf{Two tiers, and the difference
matters.} Entries carrying a full critical summary are lifted \emph{verbatim} from the
archived daily issue that first published them --- they were read, and the stated
weaknesses are assessed weaknesses. The remainder appear as precis or metadata-only
entries in the lighter style: those records were surfaced and classified but
\textbf{not read}, either because no abstract was available in Crossref, Europe PMC or
Semantic Scholar (15 records this week, 11\% of the corpus) or because they entered on
metadata alone. \textbf{Absence of criticism in the lighter tier means the record is
unassessed, not that it is sound.} Do not cite a precis entry as though it had been
evaluated.

\vspace{3mm}

\input{digest_body}

\clearpage

% ==================================================================== TABLE
\section*{\sffamily\bfseries\color{Deep}Pooled corpus register}
\vspace{-1.5mm}{\color{Teal}\rule{\linewidth}{1.1pt}}\vspace{2.5mm}

{\fontsize{6.7}{8.4}\selectfont
\setlength{\tabcolsep}{2.4pt}
\begin{longtable}{@{}>{\RaggedRight}p{25mm}>{\RaggedRight}p{41mm}>{\RaggedRight}p{27mm}%
>{\RaggedRight}p{23mm}>{\RaggedRight}p{18mm}>{\RaggedRight}p{15mm}@{}}
\toprule
\textsf{\textbf{First author}} & \textsf{\textbf{Focus}} & \textsf{\textbf{Journal}} &
\textsf{\textbf{Design}} & \textsf{\textbf{Metric}} & \textsf{\textbf{Region}}\\
\midrule
\endfirsthead
\toprule
\textsf{\textbf{First author}} & \textsf{\textbf{Focus}} & \textsf{\textbf{Journal}} &
\textsf{\textbf{Design}} & \textsf{\textbf{Metric}} & \textsf{\textbf{Region}}\\
\midrule
\endhead
\multicolumn{6}{@{}l@{}}{\cellcolor{Mist}\textsf{\textbf{\textcolor{Deep}{Neuro / mental health}}}}\\[0.6mm]
\href{https://doi.org/10.1001/jamanetworkopen.2026.25015}{Dales} & Wildfire-sourced PM$_{2.5}$ and ED visits for migraine and other primary headache & \textit{JAMA Netw Open} & Case-crossover & PM$_{2.5}$ (wildfire) & Canada\\
\href{https://doi.org/10.1093/ageing/afag215}{Glover} & Umbrella meta-review of urban environmental exposures and cognitive health & \textit{Age Ageing} & Umbrella review / meta-analysis & PM$_{2.5}$ & Global\\
\href{https://doi.org/10.1016/j.nsa.2026.107019}{Goossens} & In utero air pollution exposure and pre-/postnatal brain development & \textit{Neurosci Appl} & Narrative review & PM$_{2.5}$ / PM$_{10}$ / PM$_{\mathrm{coarse}}$ & Global\\
\href{https://doi.org/10.1016/j.isci.2026.116617}{Ma} & PM$_{2.5}$-induced neurotoxicity and depression: exposure risks and mechanisms & \textit{iScience} & Narrative review & PM$_{2.5}$ + components & Global\\
\href{https://doi.org/10.3760/cma.j.cn112150-20250804-00752}{Pan} & Maternal PM exposure and offspring neurodevelopment: effects and mechanisms & \textit{Zhonghua Yu Fang Yi Xue Za Zhi} & Narrative review & PM (bulk) & China\\
\addlinespace[0.9mm]
\multicolumn{6}{@{}l@{}}{\cellcolor{Mist}\textsf{\textbf{\textcolor{Teal}{Cardiovascular \& metabolic}}}}\\[0.6mm]
\href{https://doi.org/10.3760/cma.j.cn112338-20251120-00843}{Huang} & PM$_{2.5}$ and dyslipidaemia in rural Chinese children; dietary effect modification & \textit{Zhonghua Liu Xing Bing Xue Za Zhi} & Cross-sectional (multistage) & PM$_{2.5}$ & China\\
\href{https://doi.org/10.32394/pe/218966}{Kazmierska} & Air pollution and cardiovascular disease: review of current data & \textit{Przegl Epidemiol} & Narrative review & PM$_{2.5}$ / PM$_{10}$ & Global\\
\href{https://doi.org/10.1007/s00125-026-06814-2}{Xi} & Air pollution, built environment and temperature vs. prediabetes phenotypes & \textit{Diabetologia} & Prospective cohort & PM$_{2.5}$ & Germany\\
\href{https://doi.org/10.1161/JAHA.125.046917}{Yang} & Short-term PM$_{2.5}$ and hospitalisation/ED use in atrial fibrillation & \textit{J Am Heart Assoc} & Case-crossover & PM$_{2.5}$ (wildfire-influenced) & USA\\
\addlinespace[0.9mm]
\multicolumn{6}{@{}l@{}}{\cellcolor{Mist}\textsf{\textbf{\textcolor{Sky}{Reproductive \& developmental}}}}\\[0.6mm]
\href{https://doi.org/10.1016/j.jes.2025.10.043}{Chen} & Prenatal air pollution, fetal growth trajectory phenotypes and child neurodevelopment & \textit{J Environ Sci (China)} & Prospective cohort & PM$_{2.5}$ / PM$_{10}$ & China\\
\href{https://doi.org/10.1016/j.ecoenv.2026.120566}{Huang} & Joint prenatal bisphenol + air-pollutant mixture exposure and preschool IQ & \textit{Ecotoxicol Environ Saf} & Prospective cohort (mixtures) & PM$_{2.5}$ / PM$_{10}$ & China\\
\href{https://doi.org/10.1007/s40572-026-00552-8}{Peterson} & Ambient air pollution, PAHs and pubertal development: critical review & \textit{Curr Environ Health Rep} & Critical review & PM$_{2.5}$ components & Global\\
\href{https://doi.org/10.1016/j.envres.2026.125313}{Xie} & Long-term PM$_{2.5}$ and timing of pubertal development in boys and girls & \textit{Environ Res} & Prospective cohort & PM$_{2.5}$ & China\\
\addlinespace[0.9mm]
\multicolumn{6}{@{}l@{}}{\cellcolor{Mist}\textsf{\textbf{\textcolor{Sage}{Respiratory \& allergic}}}}\\[0.6mm]
\href{https://doi.org/10.1007/s00484-026-03289-z}{Bayar Muluk} & Climate change and otolaryngology: airway disease impacts and emission reduction & \textit{Int J Biometeorol} & Narrative review & PM$_{2.5}$ & Global\\
\href{https://doi.org/10.4103/mgr.MEDGASRES-D-25-00164}{Liao} & Short-term air pollution, respiratory symptoms and lung function in asthma & \textit{Med Gas Res} & Panel study & PM$_{2.5}$ / PM$_{10}$ & China\\
\href{https://doi.org/10.3390/epidemiologia7040097}{Sanchez de Prada} & Ambient PM and asthma-coded admissions in children 0-3 y, Spain & \textit{Epidemiologia} & Ecological panel & PM$_{2.5}$ / PM$_{10}$ & Spain\\
\href{https://doi.org/10.1016/j.jhazmat.2026.143076}{Yang} & Allergenic vs. toxic risk of airborne seed hairs compared with ambient PM$_{2.5}$ & \textit{J Hazard Mater} & Field measurement & PM$_{2.5}$ + bioaerosol & China\\
\addlinespace[0.9mm]
\multicolumn{6}{@{}l@{}}{\cellcolor{Mist}\textsf{\textbf{\textcolor{Violet}{Mechanistic toxicology}}}}\\[0.6mm]
\href{https://doi.org/10.1016/j.jes.2025.12.074}{Ding} & PM$_{2.5}$ promotes vascular calcification via miR-27b-5p/HES1 signalling & \textit{J Environ Sci (China)} & In vitro + organoid + murine & PM$_{2.5}$ & China\\
\href{https://doi.org/10.1007/s00360-026-01698-5}{Maraschi} & Osmoregulatory disruption and metal bioaccumulation from settleable atmospheric PM & \textit{J Comp Physiol B} & Controlled exposure (ecotox) & Settleable APM & Brazil\\
\href{https://doi.org/10.1016/j.jes.2025.11.018}{Wu} & Macrophage NLRP3 inflammasome in real-ambient PM-induced hepatic glucose dysregulation & \textit{J Environ Sci (China)} & Real-ambient murine + co-culture & Real-ambient PM & China\\
\addlinespace[0.9mm]
\multicolumn{6}{@{}l@{}}{\cellcolor{Mist}\textsf{\textbf{\textcolor{Clay}{Occupational \& indoor}}}}\\[0.6mm]
\href{https://doi.org/10.3389/fmedt.2026.1796828}{Chen} & Nested suction instrument to cut intraoperative surgical-smoke PM$_{2.5}$ & \textit{Front Med Technol} & Bench / device experiment & PM$_{2.5}$ (surgical smoke) & China\\
\href{https://doi.org/10.1093/annweh/wxag055}{Cole} & Determinants of respirable crystalline silica exposure in Australian tunnelling & \textit{Ann Work Expo Health} & Observational exposure study & Respirable dust / RCS & Australia\\
\href{https://doi.org/10.1016/j.jhazmat.2026.143081}{Ding} & Stage-resolved hair Zn-isotope dynamics as a firefighting combustion-exposure tracer & \textit{J Hazard Mater} & Repeated-measures biomarker & Combustion-derived PM & China\\
\href{https://doi.org/10.3390/jox16040119}{Ramsden} & Xenobiotic hazards in aircraft cabin air & \textit{J Xenobiot} & Review + risk economics & Ultrafine particles & Global\\
\addlinespace[0.9mm]
\multicolumn{6}{@{}l@{}}{\cellcolor{Mist}\textsf{\textbf{\textcolor{Amber}{Sensing, forecasting \& instrumentation}}}}\\[0.6mm]
\href{https://doi.org/10.1016/j.jenvman.2026.130515}{Shi} & Adaptive spatiotemporal graph transformer for multi-site multi-step PM$_{2.5}$ forecasting & \textit{J Environ Manage} & Deep-learning model & PM$_{2.5}$ & China / India\\
\href{https://doi.org/10.3390/bios16070380}{Tsuchiya} & Seasonal and indoor-outdoor characteristics of fluorescent aerosol particles in offices & \textit{Biosensors} & Real-time sensor campaign & Bioaerosol / coarse & Japan\\
\addlinespace[0.9mm]
\multicolumn{6}{@{}l@{}}{\cellcolor{Mist}\textsf{\textbf{\textcolor{Amber}{Exposure assessment \& modelling}}}}\\[0.6mm]
\href{https://doi.org/10.1029/2025GH001621}{Solorzano-Araque} & Multiscale exposure modelling: LOTOS-EUROS + mobility surveys + particle toxicity & \textit{GeoHealth} & CTM + mobility integration & PM$_{2.5}$ / PM$_{10}$ & Colombia\\
\addlinespace[0.9mm]
\multicolumn{6}{@{}l@{}}{\cellcolor{Mist}\textsf{\textbf{\textcolor{Coral}{Burden, policy \& mitigation}}}}\\[0.6mm]
\href{https://doi.org/10.1002/hsr2.72886}{Karamzad} & Ischaemic heart disease burden in MENA 1990-2021; PM attributable fraction 28.2\% & \textit{Health Sci Rep} & GBD secondary analysis & Ambient PM & MENA\\
\href{https://doi.org/10.1371/journal.pone.0347989}{Rees} & HEPA filtration to reduce respiratory infection in care homes & \textit{PLoS One} & RCT process evaluation & Filtered PM & UK\\
\href{https://doi.org/10.1016/j.jes.2026.01.009}{Shan} & Policy-driven emission reductions and post-pandemic rebound, high-resolution inventories & \textit{J Environ Sci (China)} & Emission inventory & PM emissions & China\\
\addlinespace[0.9mm]
\multicolumn{6}{@{}l@{}}{\cellcolor{Mist}\textsf{\textbf{\textcolor{Slate}{Other clinical endpoints}}}}\\[0.6mm]
\href{https://doi.org/10.3390/dj14070440}{Lupita} & PM$_{2.5}$, preventive-care capacity and caries/edentulism burden across the EU & \textit{Dent J} & Ecological panel & PM$_{2.5}$ & EU-27\\
\href{https://doi.org/10.1016/j.jdin.2026.06.008}{Park} & Long-term PM exposure and risk of androgenetic alopecia & \textit{JAAD Int} & Nationwide cohort & PM$_{2.5}$ / PM$_{10}$ & South Korea\\
\href{https://doi.org/10.3171/2026.3.PEDS2615}{Sun} & Environmental exposure and area deprivation vs. paediatric hydrocephalus outcomes & \textit{J Neurosurg Pediatr} & Registry cohort + geospatial & PM$_{2.5}$ (satellite-derived) & USA\\
\addlinespace[0.9mm]
\bottomrule
\end{longtable}}


\vspace{4mm}
\begin{keybox}
{\sffamily\bfseries\footnotesize\color{Deep}Provenance and method}\\[1.2mm]
{\fontsize{7.2}{8.8}\selectfont
This rollup is an \textbf{aggregation only}. It pools \texttt{state/corpus/*.json},
\texttt{state/metrics.csv} and \texttt{state/trials.json} over
\PERIODSTART{}--\PERIODEND{} via \texttt{PMRW\_START}/\texttt{PMRW\_END} and
\textbf{re-queries no literature API}, so every record was screened and deduplicated in
the daily issue that first carried it. Per-record provenance, source-query windows,
exclusions and known metadata defects are documented in those daily issues and are not
restated here. \textbf{Daily issues covered: 9, 10, 11, 12, 13, 14 and 15 August 2026 ---
seven of seven, no missed runs and no gap in entry-date coverage.}
\textbf{One exception to the no-re-query rule, made deliberately and recorded:} the 15
metadata-only DOIs were re-resolved against Crossref to test whether abstracts had since
been deposited. \textbf{Zero of fifteen} had; six authorships were recovered and written
back to the store (see ``A label defect repaired by this rollup'').
\textbf{Known defects carried into this rollup, stated rather than hidden:}
(i) six endpoint labels meaning ``no health endpoint'' are unmapped in
\texttt{ENDPOINT\_CANON}, so \texttt{f2b} understates that bar by 30 records --- the true
figure is 65 of 134; (ii) journal names are free text and Research Square appears under
two spellings, understating preprint share; (iii) 15 records carry no abstract from any
checked source and are unassessed; (iv) the 11 August PDF as shipped shows journal names
in place of six author names, corrected in the store but not in that PDF.
Effect estimates are transcribed from abstracts and use heterogeneous exposure contrasts
--- see the forest-plot caption; \textbf{no meta-analytic summary is computed and none
should be}. Trend and inflection statements are judgements over a seven-day window and are
not statistical tests. Content derived from publisher metadata; abstracts remain the
copyright of their respective publishers.\par}
\end{keybox}

\end{document}
